\vspace{-12pt}
\section{Related Work}
\vspace{-8pt}

Currently, it is a requirement for most high-quality systematic reviews to retrieve literature using a Boolean query~\cite{chandler2019cochrane,suhail2013methods}. Given that a Boolean query retrieves studies in an unordered set, it is also a requirement that all of the studies must be screened (assessed) for inclusion in the systematic review~\cite{chandler2019cochrane}. It is currently becoming more common for a ranking to be induced over this set of studies in order to begin downstream processes of the systematic review earlier~\cite{norman2019measuring}, e.g., acquiring the full-text of studies or results extraction. This ranking of studies has come to be known as `screening prioritisation', as popularised by the CLEF TAR tasks which aimed to automate these early stages of the systematic review creation pipeline~\cite{kanoulas2017clef,kanoulas2018clef,kanoulas2019clef}. As a result, in recent years there has been an uptake in Information Retrieval approaches to enable screening prioritisation~\cite{miwa2014reducing,chen2017ecnu,alharbi2017ranking,wu2018ecnu,alharbi2018retrieving,scells2017ltr,lee2018seed,lagopoulos2018learning,abualsaud2018system,zou2018technology,scells2020clf}. The vast majority of screening prioritisation use a different representation than the original Boolean query for ranking. Often, a separate query must be used to perform ranking, which may not represent the same information need as the Boolean query. Instead, the SDR method by Lee and Sun~\cite{lee2018seed} forgoes the query all together and uses studies that have a high likelihood of relevance, seed studies~\cite{suhail2013methods}, to rank the remaining studies. 
%These studies that have a high likelihood of relevance are called `seed studies', as they \textit{seed} the initial Boolean query formulation process~\cite{suhail2013methods}. 
These are studies that are known a priori to the query formulation step. The use of documents for ranking is similar to the task of query-by-document~\cite{yang2009query,lv2011learning} which has also been used extensively in domain-specific applications~\cite{kim2014diversifying,kim2014automatic,kim2015improving}. However, as Lee and Sun note, the majority of these methods try to extract key phrases or concepts from these documents to use for searching. SDR differentiates itself from these as the intuition is that the entire document is a relevance signal, rather than certain meaningful sections. Given the relatively short length of documents here (i.e., abstracts of studies), this intuition is more meaningful than other settings where the length of a document may be much longer.